% =========================================================
% Chapter 6 stakeholder cards (Resident / Public, City Planner, Police, Manager, Researcher,
% Governance). Shared macros are defined once here, and each stakeholder subsection (heading,
% short introduction and card) sits on its own page. Cards are not floated ([H]) so each stays
% on the page of its heading; the subsection heading names the stakeholder, so cards carry no
% role tag.
% Layout: every card has the same structure (\cardpair): two panels side by side, then the two
% text sections side by side under them, so both text sections start on the same line; then one
% compact row of key figures and the footer. Card text is justified. Panels are sized to fill the
% space left on the page below the heading and introduction (\cardavail, measured on each page):
% \fitcard sets the card several times, adjusting the panel height to the gap, and keeps the
% largest panel height that fits, capped by the column width; if the panels would stay small, the
% card text is set one size smaller and the fit repeated. Cards therefore print at full text
% width; only if nothing fits is a card scaled down to \cardavail as a last resort.
% All map and chart panels are drawn by the dashboard's own builders (export_dashboard_panels.py),
% so the planner card's GeoSensitivity map is the same map as on the dashboard.
% Text and tiles of every card are generated by the dashboard LLM pipeline
% (generate_card_narrative.py, master prompt geosmile-hull-v9 + card-mode addendum); prompts,
% evidence, raw output and checks are in <stakeholder>_card_llm.json.
%
% Preamble: \usepackage[most]{tcolorbox}, xcolor, graphicx, float (nothing else).
% Chapter6/figs/ needs every <stakeholder>_card_llm_content_openai.tex and the panel images
% named in each figure below.
% Note: this file is edited by hand; build_stakeholder_cards_chapter.py would overwrite it with
% the older layout of the individual <stakeholder>_card.tex files.
% =========================================================

\definecolor{cardInk}{HTML}{12171C}
\definecolor{cardMuted}{HTML}{57646D}
\definecolor{cardPink}{HTML}{FCE4EC}
\definecolor{cardWine}{HTML}{8A0033}
\definecolor{cardRule}{HTML}{D9DEE2}
\definecolor{cardFrame}{HTML}{7A868F}

% Figures and key terms the model marked: one pink highlight in the running text's own font.
\providecommand{\hlpink}[1]{\begingroup\setlength{\fboxsep}{0.6pt}\colorbox{cardPink}{\textcolor{cardWine}{#1}}\endgroup}
\providecommand{\pv}[1]{\hlpink{#1}}
\providecommand{\kt}[1]{\hlpink{#1}}
\providecommand{\kicker}[1]{{\color{cardMuted}\sffamily\small\bfseries #1}}
% Stat tile: compact figure over a short label, four to a row in one equal-height group per card.
\providecommand{\cardtilewidth}{0.2375\linewidth}
\providecommand{\tile}[2]{\begin{tcolorbox}[enhanced, nobeforeafter, width=\cardtilewidth,
  equal height group=\cardtilegroup, colback=cardPink, frame hidden, arc=1.5mm,
  left=4pt, right=2pt, top=2pt, bottom=2pt, valign=top, before upper=\raggedright]
  {\bfseries\color{cardWine}#1}\par
  {\sffamily\footnotesize\color{cardMuted}#2}\end{tcolorbox}}
\providecommand{\cardtiles}{\noindent\cardTileA\hfill\cardTileB\hfill\cardTileC\hfill\cardTileD\par}
% Panel height, set by \fitcard; each panel is as large as its column width and this height allow.
\newlength{\cardfigh}
\providecommand{\cardpanel}[2]{{\raggedright\hyphenpenalty=10000\kicker{#1}\par}\vspace{2pt}%
  {\centering\includegraphics[width=\linewidth,height=\cardfigh,keepaspectratio]{#2}\par}}
% Card text set justified, like the thesis body; extra stretch avoids overfull lines in narrow
% columns with highlighted terms.
\providecommand{\cardjust}{\leftskip=0pt\rightskip=0pt\parfillskip=0pt plus 1fil\emergencystretch=2em}
% Two panels side by side, then their two text sections in a second row, so both text sections
% start on the same line whatever the panel heights.
% #1/#2 left panel heading/image, #3/#4 right panel heading/image, #5/#6 left text heading/text,
% #7/#8 right text heading/text (text headings may be empty).
\providecommand{\cardtexthead}[1]{\if\relax\detokenize{#1}\relax\else{\raggedright\hyphenpenalty=10000\kicker{#1}\par}\vspace{2pt}\fi}
\providecommand{\cardpair}[8]{%
  \begin{minipage}[t]{0.485\linewidth}\vspace{0pt}\cardpanel{#1}{#2}\end{minipage}\hfill
  \begin{minipage}[t]{0.485\linewidth}\vspace{0pt}\cardpanel{#3}{#4}\end{minipage}\par\vspace{3pt}
  \begin{minipage}[t]{0.485\linewidth}\vspace{0pt}\cardtexthead{#5}\cardjust#6\par\end{minipage}\hfill
  \begin{minipage}[t]{0.485\linewidth}\vspace{0pt}\cardtexthead{#7}\cardjust#8\par\end{minipage}\par}
% Shared card frame and footer rule.
\tcbset{stakecard/.style={enhanced, colback=white, colframe=cardFrame, arc=3mm, boxrule=0.8pt,
  left=3mm, right=3mm, top=2.5mm, bottom=2.5mm, before skip=0pt, after skip=2pt,
  fontupper=\linespread{1}\sffamily\cardfont\color{cardInk}}}
\providecommand{\cardfont}{\normalsize}
\providecommand{\cardfoot}[1]{\par\vspace{4pt}{\color{cardRule}\rule{\linewidth}{0.4pt}}\par\vspace{1pt}%
  {\color{cardMuted}\footnotesize #1\par}}
% Space left on the page for the card, measured after each introduction; room kept for a caption
% of up to five lines and the space LaTeX puts around a non-floating ([H]) figure.
\newlength{\cardavail}
\providecommand{\cardcaptionroom}{\dimexpr5.5\baselineskip+2\intextsep+\abovecaptionskip+\belowcaptionskip\relax}
\providecommand{\measurecard}{\par\setlength{\cardavail}{\dimexpr\pagegoal-\pagetotal-\cardcaptionroom\relax}}
% \fitcard{rows}: sets \cardcontent (one tcolorbox) with trial panels, then adjusts the panel
% height by the gap to \cardavail (divided among the rows of panels stacked vertically) over a few
% passes and keeps the largest panel height whose card still fits. If that leaves the panels below
% \cardminfigh (long text in a wide thesis font), the card text is set one size smaller
% (\small, then \footnotesize) and the fit repeated, so panels stay legible and the card keeps the
% full text width; only if nothing fits is the card scaled down to \cardavail.
\newsavebox{\cardbox}
\newlength{\cardheight}
\newlength{\cardbesth}
\providecommand{\cardtrialh}{5cm}
\providecommand{\cardminfigh}{5.5cm}
\providecommand{\setcard}{\sbox{\cardbox}{\begin{minipage}{\linewidth}\cardcontent\end{minipage}}%
  \setlength{\cardheight}{\dimexpr\ht\cardbox+\dp\cardbox\relax}%
  \ifdim\cardheight>\cardavail\else\ifdim\cardfigh>\cardbesth\setlength{\cardbesth}{\cardfigh}\fi\fi}
\providecommand{\fitstep}[1]{\setlength{\cardfigh}{\dimexpr\cardfigh+(\cardavail-\cardheight)/#1\relax}%
  \ifdim\cardfigh<2.5cm\setlength{\cardfigh}{2.5cm}\fi\ifdim\cardfigh>\cardavail\setlength{\cardfigh}{\cardavail}\fi
  \setcard}
\providecommand{\fitpass}[1]{\setlength{\cardbesth}{0pt}\setlength{\cardfigh}{\cardtrialh}\setcard
  \fitstep{#1}\fitstep{#1}\fitstep{#1}\fitstep{#1}\fitstep{#1}\fitstep{#1}%
  \ifdim\cardbesth>0pt\setlength{\cardfigh}{\cardbesth}\setcard\fi}
% Each text size has its own tile height group, so tiles are only as tall as their own labels.
\providecommand{\fitcard}[1]{\edef\cardtilebase{\cardtilegroup}%
  \def\cardfont{\normalsize}\edef\cardtilegroup{\cardtilebase-n}\fitpass{#1}%
  \ifdim\cardbesth<\cardminfigh\def\cardfont{\small}\edef\cardtilegroup{\cardtilebase-s}\fitpass{#1}\fi
  \ifdim\cardbesth<\cardminfigh\def\cardfont{\footnotesize}\edef\cardtilegroup{\cardtilebase-f}\fitpass{#1}\fi
  {\centering\ifdim\cardheight>\cardavail\resizebox*{!}{\cardavail}{\usebox{\cardbox}}\else\usebox{\cardbox}\fi\par}}
% Placeholders, replaced by each card's generated content file.
\providecommand{\cardHeadline}{[headline]}
\providecommand{\cardSecA}{[section 01]}
\providecommand{\cardSecB}{[section 02]}
\providecommand{\cardFooter}{[footer]}
\providecommand{\cardModel}{[model]}
\providecommand{\cardTileA}{\tile{--}{--}}
\providecommand{\cardTileB}{\tile{--}{--}}
\providecommand{\cardTileC}{\tile{--}{--}}
\providecommand{\cardTileD}{\tile{--}{--}}


% =========================================================
\clearpage
\input{Chapter6/figs/public_card_llm_content_openai.tex}
\subsection{Resident / Public}
% =========================================================
For residents and the public, the questions are \emph{what the model pays attention to across
Hull}, and \emph{what it weighs in a place like theirs} (Figure~\ref{fig:stakeholder-public}).\par\medskip
\def\cardtilegroup{publictiles}
\measurecard

\begin{figure}[H]
\def\cardcontent{\begin{tcolorbox}[stakecard]
\setlength{\parskip}{0pt}\raggedright%
{\color{cardMuted}What the model pays attention to, across Hull and in a typical area}\par\vspace{4pt}
\cardpair{01 \textperiodcentered{} Across Hull}{Chapter6/figs/dashboard_public_short_E01012806_1.png}
  {02 \textperiodcentered{} A typical area: Stoneferry}{Chapter6/figs/dashboard_public_short_E01012806_2.png}
  {}{\cardSecA}
  {}{\cardSecB}
\vspace{5pt}
\cardtiles
\cardfoot{\cardFooter{} Panels are the dashboard's public view.}
\end{tcolorbox}}
\fitcard{1}
\caption{Resident / Public card: where predictions depend most on knowing an area's location, what the
model weighs in a typical area (Stoneferry), and the LLM-generated plain-language explanation.}
\label{fig:stakeholder-public}
\end{figure}


% =========================================================
\clearpage
\input{Chapter6/figs/planner_card_llm_content_openai.tex}
\subsection{City Planner / Local Authority}
% =====================================================================
For a city planner, two questions are most informative: \emph{which areas matter most to the
model}, and \emph{which built-environment features it relies on, and where}
(Figure~\ref{fig:stakeholder-planner}).\par\medskip
\def\cardtilegroup{plannertiles}
\measurecard

\begin{figure}[H]
\def\cardcontent{\begin{tcolorbox}[stakecard]
\setlength{\parskip}{0pt}\raggedright%
{\color{cardMuted}Important areas and built-environment reliance}\par\vspace{4pt}
\cardpair{01 \textperiodcentered{} Most important areas}{Chapter6/figs/dashboard_planner_short_E01033104_1.png}
  {02 \textperiodcentered{} Built-environment share}{Chapter6/figs/dashboard_planner_short_E01033104_2.png}
  {}{\cardSecA}
  {}{\cardSecB}
\vspace{5pt}
\cardtiles
\cardfoot{\cardFooter{} Map 01: signed GeoSensitivity (sign is not a crime direction). Map 02:
share of local importance, not an effect on crime.}
\end{tcolorbox}}
\fitcard{1}
\caption{City Planner card: GeoSensitivity (left) and built-environment share of local importance
(right); text generated by the LLM pipeline.}
\label{fig:stakeholder-planner}
\end{figure}


% =========================================================
\clearpage
\input{Chapter6/figs/police_card_llm_content_openai.tex}
\subsection{Police / Community Safety}
% =========================================================
For police and community safety teams, two questions are most informative: \emph{whether the
model reads an area through its recent months or its long-run crime level}, and \emph{where
recorded crime is rising or falling} (Figure~\ref{fig:stakeholder-police}).\par\medskip
\def\cardtilegroup{policetiles}
\measurecard

\begin{figure}[H]
\def\cardcontent{\begin{tcolorbox}[stakecard]
\setlength{\parskip}{0pt}\raggedright%
{\color{cardMuted}Sensitivity, recent or long-run signals and recorded change}\par\vspace{4pt}
\cardpair{01 \textperiodcentered{} Recent or long-run signal}{Chapter6/figs/dashboard_police_short_E01033104_1.png}
  {02 \textperiodcentered{} Recorded change}{Chapter6/figs/dashboard_police_short_E01033104_2.png}
  {}{\cardSecA}
  {}{\cardSecB}
\vspace{5pt}
\cardtiles
\cardfoot{\cardFooter{} Panels are the dashboard's police view for Old Town (029E).}
\end{tcolorbox}}
\fitcard{1}
\caption{Police / Community Safety card: the dashboard's police panels for Old Town, with the
LLM-generated narrative.}
\label{fig:stakeholder-police}
\end{figure}


% =========================================================
\clearpage
\input{Chapter6/figs/manager_card_llm_content_openai.tex}
\subsection{Manager / Strategic Decision Maker}
% =========================================================
For a manager, for example in local government or insurance, the questions combine both
branches: \emph{what drives the model and where it is most sensitive}, and \emph{what can be set
aside, and how far the results can be trusted}
(Figure~\ref{fig:stakeholder-manager}).\par\medskip
\def\cardtilegroup{managertiles}
\measurecard

\begin{figure}[H]
\def\cardcontent{\begin{tcolorbox}[stakecard]
\setlength{\parskip}{0pt}\raggedright%
{\color{cardMuted}What drives the model, where it is sensitive, and what can be set aside}\par\vspace{4pt}
\cardpair{01 \textperiodcentered{} Most sensitive areas}{Chapter6/figs/dashboard_manager_short_E01033104_1.png}
  {02 \textperiodcentered{} What drives the model}{Chapter6/figs/dashboard_manager_short_E01033104_2.png}
  {Drivers and exceptions}{\cardSecA}
  {What can be set aside, and confidence}{\cardSecB}
\vspace{5pt}
\cardtiles
\cardfoot{\cardFooter{} Map 01: Geo branch. Panel 02: Feature branch, share of local importance,
not an effect on crime.}
\end{tcolorbox}}
\fitcard{1}
\caption{Manager / Strategic Decision Maker card: where the model is most sensitive (Geo branch)
and what drives it (Feature branch, Hull against Bilton Grange), with the LLM-generated synthesis.}
\label{fig:stakeholder-manager}
\end{figure}


% =========================================================
\clearpage
\input{Chapter6/figs/researcher_card_llm_content_openai.tex}
\subsection{GeoAI Researcher / Model Analyst}
% =========================================================
For a GeoAI researcher, the questions are technical: \emph{how Geo-SMILE behaves across Hull},
and \emph{how faithful, spatially coherent and reproducible its explanations are}
(Figure~\ref{fig:stakeholder-researcher}).\par\medskip
\def\cardtilegroup{researchertiles}
\measurecard

\begin{figure}[H]
\def\cardcontent{\begin{tcolorbox}[stakecard]
\setlength{\parskip}{0pt}\raggedright%
{\color{cardMuted}How Geo-SMILE behaves across Hull, and how far its explanations can be trusted}\par\vspace{4pt}
\cardpair{01 \textperiodcentered{} Surrogate against exact masking}{Chapter6/figs/dashboard_researcher_short_hull_1.png}
  {02 \textperiodcentered{} Top feature, local effect}{Chapter6/figs/dashboard_researcher_short_hull_2.png}
  {Explanation behaviour across Hull}{\cardSecA}
  {Fidelity, intervention and reproducibility}{\cardSecB}
\vspace{5pt}
\cardtiles
\cardfoot{\cardFooter{} Panels are the dashboard's Hull-wide researcher view.}
\end{tcolorbox}}
\fitcard{1}
\caption{GeoAI Researcher / Model Analyst card: Hull-wide surrogate check (Geo branch), local
effect of the most important feature (Feature branch), explanation-quality figures and the
LLM-generated technical narrative.}
\label{fig:stakeholder-researcher}
\end{figure}


% =========================================================
\clearpage
\input{Chapter6/figs/governance_card_llm_content_openai.tex}
\subsection{Governance / Assurance}
% =========================================================
For governance and assurance, the questions are \emph{whether the explanation is traceable and
reproducible}, and \emph{which limitations require explicit review before the evidence is
relied upon} (Figure~\ref{fig:stakeholder-governance}).\par\medskip
\def\cardtilegroup{governancetiles}
\measurecard

\begin{figure}[H]
% Only the two seed-reproducibility tiles; the rank-ratio and ATT Faithfulness tiles are omitted.
\def\cardtilewidth{0.49\linewidth}
\def\cardtiles{\noindent\cardTileA\hfill\cardTileB\par}
\def\cardcontent{\begin{tcolorbox}[stakecard]
\setlength{\parskip}{0pt}\raggedright%
{\color{cardMuted}Provenance, reproducibility and the limitations to review before reliance}\par\vspace{4pt}
\cardpair{01 \textperiodcentered{} Reproducibility across seeds}{Chapter6/figs/dashboard_governance_short_hull_1.png}
  {02 \textperiodcentered{} Stability under neighbour additions}{Chapter6/figs/dashboard_governance_short_hull_2.png}
  {Provenance and reproducibility}{\cardSecA}
  {Limitations requiring review}{\cardSecB}
\vspace{5pt}
\cardtiles
\cardfoot{\cardFooter{} Panels are the dashboard's governance view.}
\end{tcolorbox}}
\fitcard{1}
\caption{Governance / Assurance card: seed reproducibility and neighbour-addition stability of
the Geo ranking, key assurance figures and the LLM-generated review of provenance and
limitations.}
\label{fig:stakeholder-governance}
\end{figure}
